\documentclass[11pt]{article}
\usepackage[T1]{fontenc}\usepackage{lmodern}
\usepackage[margin=1in]{geometry}
\usepackage{amsmath,amssymb,amsthm,microtype}
\usepackage[hidelinks]{hyperref}
\newtheorem{theorem}{Theorem}[section]
\newtheorem{lemma}[theorem]{Lemma}
\newtheorem{proposition}[theorem]{Proposition}
\newcommand{\R}{\mathbb R}\newcommand{\E}{\mathbf E}
\newcommand{\SU}{\operatorname{SU}(2)}
\newcommand{\KL}{\operatorname{D}}
\newcommand{\lean}[2]{\par\smallskip\noindent\textit{Source correspondence:}
\texttt{\detokenize{#1}}\par\noindent\texttt{\detokenize{#2}}\par\smallskip}
\title{Conditional resampling and retained boundary entropy\\in finite SU(2) models}
\author{Project Stardust}\date{October 29, 2026}
\begin{document}\maketitle
\begin{abstract}
A conditional update that retains its boundary can remove all error in the
conditional law without removing any error in the boundary marginal. We give
complete probability-level proofs of the update law, idempotence, stationary
family and exact finite relative-entropy decomposition. For bounded finite-word
SU(2) interactions we construct the globally normalized joint law and its
actual conditional kernel, including the interaction-induced boundary marginal.
Conditioning that continuous boundary on a nontrivial open hemisphere gives a
stationary perturbation with finite strictly positive entropy against the
original joint law. Every strict global entropy contraction factor therefore
fails for this update, at every fixed finite coupling. The obstruction uses
the original continuous-boundary reference, in addition to an elementary
two-atom example. A separate state-replacement loop kernel has genuine fixed-coupling entropy
decay, with finite entropy after one step even from a point mass. A zero-trace
band supplies a genuine two-step Haar reset with explicit pair-count bounds;
an exact zero-trace row separates this mechanism from Wilson transport. Classical
conditional-expectation and heat-bath mechanisms
are credited; no obstruction to dynamics updating all coordinates or to a
physical Yang--Mills mass gap is asserted.
\end{abstract}

\section{The distinction to be proved}
Fix a boundary variable $p$ and an interior variable $x$. A heat-bath step
retains $p$ and replaces $x$ by a draw from a specified conditional law
$K_p$. Thus it enforces the desired conditional distribution in one step.
It does not force the boundary marginal to equal any chosen reference
marginal. This elementary fact becomes useful when an entropy chain rule is
mistaken for a strictly positive global dissipation estimate.

The frozen verified source for this exposition is commit
\begin{quote}\small\texttt{bab688bb6140cd6316c80680095a914ea0d346db}.\end{quote}
The accompanying map supplies exact declaration headers, hashes and links.
The main construction is already in that library; no new research proof,
mathematical source change or external submission is made here.
We work with probability measures on standard Borel spaces, sufficient for
all continuous compact-group examples. The formal generic kernel results
also cover more general measurable spaces and, in their stated forms,
finite or sigma-finite measures. We do not identify those broader hypotheses
with the narrower probability-level exposition.

For probability measures $P,Q$, relative entropy means
\[
 \KL(P\Vert Q)=\int f\log f\,dQ\quad\text{if }P=fQ,
 \qquad \KL(P\Vert Q)=+\infty\quad\text{if }P\not\ll Q.
\]
Use $0\log0=0$. This is an extended nonnegative quantity. The integral is
well defined because the negative part of $f\log f$ is bounded by $1/e$.
For the exact decomposition below we assume finite input entropy, so all
terms can be handled without subtracting infinities. The actual SU(2)
counterexample will have entropy equal to a specified finite positive number.

\section{An actual boundary-preserving kernel}
Let $\mathsf P,\mathsf E$ be standard Borel spaces, and let
$K:\mathsf P\to\mathcal P(\mathsf E)$ be a Markov kernel. Write
$\pi\otimes K$ for the joint measure
\[
 (\pi\otimes K)(B)=\int\!\int1_B(p,x)\,K_p(dx)\,\pi(dp).
\]
Define a kernel $T$ on $\mathsf P\times\mathsf E$ by
\begin{equation}\label{eq:T}
 T((p,x),B)=\int1_B(p,y)\,K_p(dy).
\end{equation}
For a law $J$, denote its boundary marginal by $J_{\mathsf P}$ and its
updated law by $JT$. This notation corresponds to kernel composition in the
formal sources, whose composition symbol is written in the opposite order.

\begin{theorem}[Update law and stationary family]\label{thm:update}
The kernel $T$ is Markov, and for every probability law $J$,
\begin{equation}\label{eq:update}
 JT=J_{\mathsf P}\otimes K.
\end{equation}
It preserves the boundary marginal and satisfies $T^2=T$. Every law
$\pi\otimes K$ is stationary. Conversely, $JT=J$ exactly when
$J=J_{\mathsf P}\otimes K$, equivalently when $K$ is a conditional kernel
for $J$.
\end{theorem}
\begin{proof}
For each input $(p,x)$, \eqref{eq:T} is the pushforward of the probability
$K_p$ under $y\mapsto(p,y)$. For each measurable output set $B$, the kernel
integration theorem makes its value measurable in $p$, and hence in $(p,x)$.
This proves that it is Markov. On rectangles $A\times C$, Tonelli gives
\[
 (JT)(A\times C)=\int1_A(p)K_p(C)\,J(d p,d x)
 =\int_A K_p(C)\,J_{\mathsf P}(dp).
\]
Rectangles determine probability measures on the product sigma algebra,
so \eqref{eq:update} follows. Taking $C=\mathsf E$ proves marginal
preservation. Applying the same formula twice gives
$(JT)T=(J_{\mathsf P}\otimes K)T=J_{\mathsf P}\otimes K$.
The same rectangle calculation for a point input proves the kernel identity
$T^2=T$. The stationary statements follow directly from \eqref{eq:update}
and the definition of a conditional kernel.
\end{proof}
\lean{ConditionalResamplingKernel}{update_markov; update_law; boundary_preserved; update_idempotent}
\lean{ConditionalResamplingKernel}{every_boundary_stationary; invariant_iff_conditional}
Different boundary marginals yield different stationary laws: the first
marginal of $\pi\otimes K$ is $\pi$. Thus stationary nonuniqueness here is
not ambiguity about versions of a conditional distribution. It is a
consequence of a coordinate that the update retains exactly.

\section{Exact finite-entropy accounting}
Fix a reference boundary probability $\tau$ and set $Q=\tau\otimes K$.
By the preceding theorem, $Q$ is stationary. The following proof explicitly
identifies the density after one update.

\begin{theorem}[Marginal and conditional entropy]\label{thm:entropy}
Suppose $\KL(J\Vert Q)<\infty$. Then
\begin{align}
 \KL(JT\Vert Q)&=\KL(J_{\mathsf P}\Vert\tau),\label{eq:after}\\
 \KL(J\Vert Q)&=\KL(JT\Vert Q)+\KL(J\Vert JT).\label{eq:chain}
\end{align}
All three quantities in \eqref{eq:chain} are finite and nonnegative.
\end{theorem}
\begin{proof}
Finite input entropy implies $J=fQ$ for a nonnegative density $f$ with
$\int f\,dQ=1$. Define
\[
 m(p)=\int f(p,x)\,K_p(dx).
\]
Tonelli gives $J_{\mathsf P}=m\tau$ and hence, by \eqref{eq:update},
$JT=mQ$. Conditional Jensen for the convex function $u\log u$ gives
\[
 \int m\log m\,d\tau\leq\int f\log f\,dQ<\infty.
\]
Both negative parts are integrable, so the two finite integrals are
unambiguous. Moreover
\[
 \int |\log m|\,dJ=\int m|\log m|\,d\tau<\infty,
\]
where expressions on $m=0$ can be assigned any value because $J$ gives that
set zero mass. The density of $JT$ relative to $Q$ is $m(p)$, so its entropy
is precisely $\int m\log m\,d\tau$, proving \eqref{eq:after}.

On $m(p)>0$, put $r(p,x)=f(p,x)/m(p)$. On $m=0$, nonnegativity and the
definition of $m$ imply $f=0$ for $K_p$-almost every $x$, for $\tau$-almost
every such $p$. Thus $J\ll JT$ and $r$ is a density of $J$ relative to
$JT$. Since $\log f$ is integrable under $J$ by finite entropy, and
$\log m$ is integrable under $J$ by the preceding calculation, their
difference $\log r$ is integrable under $J$. Therefore
\[
 \KL(J\Vert Q)=\int\log f\,dJ
 =\int\log m\,dJ+\int\log r\,dJ
 =\KL(JT\Vert Q)+\KL(J\Vert JT).
\]
Nonnegativity follows from Jensen, or from $u\log u\geq u-1$ integrated
against a probability reference. The established integrability gives
finiteness of the final term as well.
\end{proof}
\lean{ConditionalResamplingEntropy}{entropy_after_update; entropy_decomposition_of_conditional}
The mapped source proves an extended-divergence chain rule using the original
measure-theoretic KL library. The displayed human proof deliberately proves
its finite-entropy probability restriction, which covers the counterexamples
and avoids any assertion that $+\infty-(+\infty)$ is defined.

There is a useful direct interpretation of the second term. For $m(p)>0$,
the input conditional law has density $f(p,x)/m(p)$ relative to $K_p$.
Consequently
\[
 \KL(J\Vert JT)=\int\KL(J(dx\mid p)\Vert K_p)\,J_{\mathsf P}(dp).
\]
This is the conditional error that the step removes. The remaining entropy
depends only on the retained boundary. On the restricted class
$J_{\mathsf P}=\tau$, one step does send every input to $Q$. That restricted
statement does not imply strict global contraction for arbitrary inputs.

\section{Stationary perturbations with positive finite entropy}
The elementary mechanism can be realized without any singular or infinite
entropy comparison.

\begin{lemma}[Conditioning a boundary event]\label{lem:event}
Let $A\subseteq\mathsf P$ be measurable with $a=\tau(A)\in(0,1)$. Define
\[
 \pi=\tau(\,\cdot\mid A)=a^{-1}1_A\tau,\qquad J_A=\pi\otimes K.
\]
Then $J_AT=J_A$ and
\begin{equation}\label{eq:evententropy}
 \KL(J_A\Vert Q)=\log(1/a)\in(0,\infty).
\end{equation}
\end{lemma}
\begin{proof}
Stationarity follows from Theorem~\ref{thm:update}. The density of $J_A$
relative to $Q$ is $a^{-1}1_A(p)$, because integrating on rectangles gives
the same joint measure. Its logarithm equals $\log(1/a)$ on the support of
$J_A$, and $J_A$ has mass one. This proves \eqref{eq:evententropy} directly.
The assumptions $0<a<1$ give both finiteness and strict positivity.
\end{proof}
\lean{ConditioningRelativeEntropy}{entropy_cond_finite; entropy_cond_exact; entropy_cond_positive}

\begin{theorem}[No strict global fraction]\label{thm:nocontraction}
Under the assumptions of Lemma~\ref{lem:event}, no $q<1$ can satisfy
\[
 \KL(JT\Vert Q)\leq q\KL(J\Vert Q)
\]
for all finite-entropy probability inputs $J$. Likewise there is no $c>0$
such that $c\KL(J\Vert Q)\leq\KL(J\Vert JT)$ for all such inputs.
\end{theorem}
\begin{proof}
Substitute $J_A$. The first inequality would say $d\leq qd$ for the finite
number $d=\log(1/a)>0$, contradicting $q<1$. The second would say
$cd\leq0$, since $J_AT=J_A$, contradicting $c,d>0$.
\end{proof}
\lean{ConditionalResamplingEntropy}{no_positive_dissipation_fraction; binary_no_strict_contraction}
For a concrete finite example let $\mathsf P=\{0,1\}$,
$\tau=\frac12(\delta_0+\delta_1)$ and $\pi=\delta_0$.
For \emph{any} Markov kernel $K$, including a strictly positive one on a
finite interior space, $\pi\otimes K$ is stationary and its entropy relative
to $\tau\otimes K$ is exactly $\log2$. This example uses a specifically
chosen two-atom boundary reference. It does not by itself prove the same
statement against a prescribed continuous interaction law; that is the task
of the next sections.

\section{Globally normalized bounded interaction laws}
Let $\rho$ and $\nu$ be probability measures on $\mathsf P$ and $\mathsf E$,
and let $V:\mathsf P\times\mathsf E\to\R$ be measurable with
$|V(p,x)|\leq B<\infty$. Define the section and total partitions
\[
 z(p)=\int e^{V(p,x)}\,\nu(dx),\qquad
 \mathcal Z=\int z(p)\,\rho(dp).
\]
Both satisfy $e^{-B}\leq z(p),\mathcal Z\leq e^B$ and are strictly positive.
The section partition is measurable by kernel integration. Set
\begin{align}
 K_p(dx)&=z(p)^{-1}e^{V(p,x)}\nu(dx),\label{eq:K}\\
 \tau(dp)&=\mathcal Z^{-1}z(p)\rho(dp),\label{eq:tau}\\
 Q(dp,dx)&=\mathcal Z^{-1}e^{V(p,x)}\rho(dp)\nu(dx).\label{eq:Q}
\end{align}

\begin{proposition}[Factorization and boundary equivalence]\label{prop:factor}
The formulas define a Markov kernel $K$ and probabilities $\tau,Q$.
Moreover $Q=\tau\otimes K$, $Q_{\mathsf P}=\tau$, and $\tau$ and $\rho$
have exactly the same null sets. In fact
\[
 e^{-2B}\leq\frac{d\tau}{d\rho}(p)\leq e^{2B}.
\]
\end{proposition}
\begin{proof}
The section density in \eqref{eq:K} is jointly measurable, nonnegative and
has integral one. It defines a Markov kernel by the kernel integration
theorem. Fubini gives the masses of \eqref{eq:tau} and \eqref{eq:Q} as one.
Multiplication of the boundary density and conditional density cancels
$z(p)$, giving the joint density \eqref{eq:Q}; alternatively integrate on
rectangles. The first marginal follows by integrating $K_p(\mathsf E)=1$.
The density bound follows from the two partition bounds. A density bounded
strictly above zero and finite has precisely the same null sets as its
reference measure.
\end{proof}
\lean{ExponentialJointDisintegration}{boundary_probability; factorization; joint_probability; conditional_kernel}
\lean{ExponentialBoundaryEquivalence}{reverse_ac; null_iff}
The original generic equivalence theorem only requires positive integrable
exponential sections and integrability of the joint exponential; the bounded
case here supplies these hypotheses directly. The explicit uniform density
bound is a human consequence of boundedness, not an asserted additional
formal declaration.

The reference boundary is $\tau$, rather than $\rho$, in the actual
globally normalized model. Unless the section partition $z(p)$ is constant,
$\rho\otimes K$ differs from $Q$. Omitting this reweighting would replace
the intended reference law and invalidate an entropy comparison claimed to
use the globally normalized interaction.

\section{Finite SU(2) words and their actual link conditional}
Let $G=\SU$, with normalized Haar probability $m$. Its matrices are unitary
of determinant one. For a finite boundary index set $I$, put
$\mathsf P=G^I$, $\mathsf E=G$, $\rho=m^{\otimes I}$ and $\nu=m$.
Choose a finite collection of words $W_j(U)$ in boundary letters $U_i$ and
their inverses. Empty words equal the identity. Matrix multiplication and
inversion are continuous on $G$, so each word is continuous, and in
particular measurable. Let the collection have $N$ words and define
\begin{equation}\label{eq:potential}
 V_\beta(U,A)=\frac\beta2\operatorname{Re}\operatorname{tr}
                   \left(A\sum_{j=1}^N W_j(U)\right).
\end{equation}
All coupling constants $\beta\in\R$ are allowed.

\begin{lemma}[Uniform finite-word bound]\label{lem:bound}
For every boundary $U$ and link $A$,
$|V_\beta(U,A)|\leq |\beta|N$.
\end{lemma}
\begin{proof}
Each product $AW_j(U)$ is unitary. Every diagonal entry of a unitary matrix
has modulus at most one, since each column has Euclidean norm one.
Therefore $|\operatorname{Re}\operatorname{tr}(AW_j(U))|\leq2$.
Linearity of the trace and the triangle inequality give a total bound
$\frac{|\beta|}{2}\sum_j2=|\beta|N$.
\end{proof}
\lean{SU2BoundaryWordKernel}{word_measurable; family_measurable; kernel_markov}
\lean{SU2StapleKernel}{exponent_bound; exponential_integrable}
The source proves the same bound through its quaternion staple radius;
the two-line matrix argument is a human reorganization of that estimate.

Apply Section~5 with $B=|\beta|N$. Denote the resulting actual joint law,
boundary marginal and conditional kernel by $Q_\beta,\tau_\beta,K_\beta$.
Their explicit densities are
\[
 Q_\beta(dU,dA)=\frac{e^{V_\beta(U,A)}}{\mathcal Z_\beta}\rho(dU)m(dA),
 \qquad
 (K_\beta)_U(dA)=\frac{e^{V_\beta(U,A)}}{z_\beta(U)}m(dA).
\]
In particular,
\begin{equation}\label{eq:actualfactor}
 Q_\beta=\tau_\beta\otimes K_\beta,
 \qquad \tau_\beta(dU)=\frac{z_\beta(U)}{\mathcal Z_\beta}\rho(dU).
\end{equation}
No formal ``conditional density'' is being used without a joint law:
these formulas construct both the joint probability and its disintegration.
The rectangle identity
\[
 Q_\beta(B\times C)=\int_B (K_\beta)_U(C)\,\tau_\beta(dU)
\]
checks the actual conditional relationship. The heat-bath update
$T_\beta((U,A),\cdot)=\delta_U\otimes(K_\beta)_U$ is Markov, idempotent,
boundary preserving, and leaves $Q_\beta$ invariant, by Theorem~\ref{thm:update}.
\lean{SU2BoundaryJointLaw}{factorization; boundary_marginal; conditional_rectangles; conditional_kernel}
\lean{SU2ConditionalHeatBath}{update_markov; update_idempotent; boundary_preserved; joint_invariant}

For every finite-entropy probability input $J$ against this actual joint
reference, Theorem~\ref{thm:entropy} now gives
\begin{align*}
 \KL(JT_\beta\Vert Q_\beta)&=\KL(J_{\mathsf P}\Vert\tau_\beta),\\
 \KL(J\Vert Q_\beta)&=\KL(JT_\beta\Vert Q_\beta)
                         +\KL(J\Vert JT_\beta).
\end{align*}
The formal source retains the extended-divergence versions as well.
\lean{SU2HeatBathEntropy}{actual_joint_entropy_after_update; actual_joint_entropy_decomposition}

\section{A stationary perturbation of the continuous-boundary reference}
We now use one genuine continuous boundary coordinate, so $I$ has one
element and $\mathsf P=G$. Every finite word collection and finite real
coupling in \eqref{eq:potential} is still permitted. Define the open event
\[
 H=\{U\in G:\operatorname{Re}U_{11}>0\}.
\]
Its complement has a nonempty open subset
$H_- =\{U:\operatorname{Re}U_{11}<0\}$.

\begin{lemma}[A nontrivial event for the actual boundary]\label{lem:hemisphere}
For every finite word interaction and $\beta\in\R$,
\[
 0<a_\beta:=\tau_\beta(H)<1.
\]
\end{lemma}
\begin{proof}
The matrix entry is continuous. The identity belongs to $H$, and $-1$
belongs to $H_-$, so both are nonempty open sets. Every nonempty open set
in a compact group has positive Haar measure: if $O$ is nonempty, its left
translates cover the group, and compactness extracts a finite subcover.
If $m(O)=0$, translation invariance would make the whole group have measure
zero, contradicting $m(G)=1$.
Thus $m(H)>0$ and $m(H^c)\geq m(H_-)>0$.
Proposition~\ref{prop:factor} says that $\tau_\beta$ has exactly the Haar
null sets, so the same strict inequalities hold for $\tau_\beta$.
Since it is a probability and $H$ is measurable, they imply
$0<\tau_\beta(H)<1$.
\end{proof}
\lean{SU2BoundaryConditionedEntropy}{haar_hemisphere_positive; haar_complement_positive}
\lean{SU2BoundaryConditionedEntropy}{boundary_null_iff; boundary_event_nontrivial}
No symmetry of the interaction has been assumed, and $a_\beta=1/2$ is
not claimed. The argument proves the strict range actually
required for finite positive entropy.

\begin{theorem}[Continuous-boundary entropy obstruction]\label{thm:su2}
Define
\[
 \widehat\tau_\beta=\tau_\beta(\,\cdot\mid H),\qquad
 \widehat Q_\beta=\widehat\tau_\beta\otimes K_\beta.
\]
These are probability measures, $\widehat Q_\beta$ is stationary for
$T_\beta$, and
\begin{equation}\label{eq:exact}
 0<\KL(\widehat Q_\beta\Vert Q_\beta)
   =\log(1/a_\beta)<\infty.
\end{equation}
For every $q<1$ the strict global contraction inequality
\[
 \KL(JT_\beta\Vert Q_\beta)\leq q\KL(J\Vert Q_\beta)
\]
fails for the actual finite-entropy input $J=\widehat Q_\beta$.
\end{theorem}
\begin{proof}
Lemma~\ref{lem:hemisphere} makes conditioning well defined.
Lemma~\ref{lem:event}, applied to $\tau_\beta,K_\beta,H$, gives stationarity
and \eqref{eq:exact}. Substitution into the proposed inequality gives
$\log(1/a_\beta)\leq q\log(1/a_\beta)$, impossible for $q<1$.
\end{proof}
\lean{SU2BoundaryConditionedEntropy}{perturbed_stationary; entropy_finite; entropy_exact}
\lean{SU2BoundaryConditionedEntropy}{entropy_positive; no_strict_contraction}
This test changes the boundary distribution absolutely continuously and
retains the correct link conditional at every retained boundary.
It is therefore stronger for this particular reference-law question than
changing the reference to a two-atom boundary law. The two examples share the
same classical mechanism, but they address different reference measures.

\section{A replacement kernel with genuine entropy decay}
For comparison we now change the state space and the update. The state is
one loop matrix $C\in G$. Put $h(C)=\frac12\operatorname{Re}\operatorname{tr}C$,
so $|h(C)|\leq1$. For $r\in\R$, define
\[
 z_r(C)=\int_G e^{r h(C)h(D)}m(dD),\qquad
 Z_r=\int_Gz_r(C)m(dC),
\]
\[
 k_r(C,D)=\frac{e^{r h(C)h(D)}}{z_r(C)},\quad
 P_r(C,dD)=k_r(C,D)m(dD),
\]
\[
 \pi_r(dC)=p_r(C)m(dC),\quad p_r(C)=\frac{z_r(C)}{Z_r}.
\]
This transition outputs $D$ as the next state; it does not output the pair
$(C,D)$ while retaining $C$. It is therefore different from the update $T$
on a boundary--interior pair. All section and total partitions are between
$e^{-|r|}$ and $e^{|r|}$, so the formulas give actual probabilities and a
Markov kernel by the same measurable-density argument as before.

\begin{proposition}[Stationary comparison bounds]\label{prop:loop}
The law $\pi_r$ is stationary for $P_r$, and
\[
 e^{-2|r|}\leq k_r(C,D),p_r(D)\leq e^{2|r|}.
\]
In particular, with $\varepsilon_r=\frac12e^{-4|r|}\in(0,1)$,
\begin{equation}\label{eq:sandwich}
 \varepsilon_r\pi_r\leq P_r(C,\cdot)\leq e^{4|r|}\pi_r
 \qquad\text{for every }C.
\end{equation}
\end{proposition}
\begin{proof}
The joint density $p_r(C)k_r(C,D)$ equals
$Z_r^{-1}e^{rh(C)h(D)}$. It is symmetric in $C,D$. Its first marginal is
$\pi_r$, so symmetry makes its second marginal $\pi_r$ as well. This is
stationarity. The two density bounds follow by dividing numerator and
partition bounds. Consequently
$k_r(C,D)\geq e^{-2|r|}\geq e^{-4|r|}p_r(D)$ and
$k_r(C,D)\leq e^{2|r|}\leq e^{4|r|}p_r(D)$.
Integrate these inequalities against $m$ and weaken the first coefficient
by a factor two to obtain \eqref{eq:sandwich}. The factor two ensures
$\varepsilon_r<1$ even at $r=0$ and matches the source's residual kernel.
\end{proof}
\lean{SU2LoopEquilibrium}{detailed_balance_density; equilibriumLaw_stationary}
\lean{SU2StationaryMinorization}{stationary_density_minorization; loopKernel_minorizes_equilibrium}
\lean{SU2PointEntropy}{reverse_density_bound; loopKernel_dominated_by_equilibrium}

We give the entropy implication with enough detail to specify the actual
residual kernel and the role of finite entropy. The following argument is
classical. Suppose $P$ is a Markov kernel on a standard Borel space, $\pi$
is stationary and $P(x,\cdot)\geq\varepsilon\pi$ for every $x$, with
$0<\varepsilon<1$. Define
\[
 R(x,\cdot)=\frac{P(x,\cdot)-\varepsilon\pi}{1-\varepsilon}.
\]
This is a Markov kernel: subtraction leaves a nonnegative measure by
minorization, its mass is $1-\varepsilon$, and measurable evaluations follow
from those of $P$. Integrating the identity
$P=\varepsilon\pi+(1-\varepsilon)R$ against $\pi$ and using stationarity
gives $\pi R=\pi$.

\begin{lemma}[Residual entropy contraction]\label{lem:residual}
For every probability $\mu$ with finite $\KL(\mu\Vert\pi)$,
\[
 \KL(\mu P\Vert\pi)\leq(1-\varepsilon)\KL(\mu\Vert\pi).
\]
\end{lemma}
\begin{proof}
We first verify data processing for $R$ in this finite-entropy setting.
Write $\mu=f\pi$. The stationary joint law $\pi(dx)R(x,dy)$ has second
marginal $\pi(dy)$. Its disintegration given $y$ supplies a reverse kernel
$\widetilde R(y,dx)$. The density of $\mu R$ relative to $\pi$ is
$g(y)=\int f(x)\widetilde R(y,dx)$, as checked by integrating against
bounded measurable functions of $y$. Jensen for $u\log u$ gives
\[
 \int g\log g\,d\pi\leq\int\!\int f(x)\log f(x)
           \widetilde R(y,dx)\pi(dy)=\int f\log f\,d\pi.
\]
Integrable negative parts justify these extended integrals; the upper bound
is finite, so $\KL(\mu R\Vert\pi)$ is finite too.

The density of $\mu P$ is $\varepsilon+(1-\varepsilon)g$.
Convexity of $u\log u$ at the mixture of $1$ and $g$ gives pointwise
\[
 (\varepsilon+(1-\varepsilon)g)
 \log(\varepsilon+(1-\varepsilon)g)
 \leq(1-\varepsilon)g\log g.
\]
Integration and the data-processing bound prove the claim.
\end{proof}
\lean{SU2ResidualEntropy}{residualKernel_markov; residualKernel_split; residualKernel_invariant}
\lean{SU2ResidualEntropy}{entropy_contraction; entropy_geometric}
The mapped formal theorem also allows infinite input entropy using extended
KL inequalities. The finite-input proof above suffices after the following
actual smoothing step and does not rely on multiplying an infinite entropy
by a geometric factor and then calling it a convergence estimate.

\begin{theorem}[Finite entropy after one step and subsequent decay]\label{thm:decay}
For every $r\in\R$ and \emph{every} probability input $\mu$ on $G$,
including a point mass,
\[
 \KL(\mu P_r\Vert\pi_r)\leq4|r|<\infty,
\]
and for all integers $n\geq0$,
\begin{equation}\label{eq:decay}
 \KL(\mu P_r^{n+1}\Vert\pi_r)
 \leq(1-\varepsilon_r)^n\KL(\mu P_r\Vert\pi_r)
 \leq4|r|(1-\varepsilon_r)^n\longrightarrow0.
\end{equation}
\end{theorem}
\begin{proof}
Integrating the upper bound in \eqref{eq:sandwich} against $\mu$ gives
$\mu P_r\leq M\pi_r$, $M=e^{4|r|}\geq1$. Its density $f$ is at most $M$.
Since $f\log f\leq f\log M$, its entropy is at most $\log M=4|r|$;
the negative part is bounded, so this also proves finiteness. Now apply
Lemma~\ref{lem:residual} successively to the finite-entropy input $\mu P_r$.
Every iterate has finite entropy by induction, giving the first inequality
in \eqref{eq:decay}. The geometric factor tends to zero because
$0<\varepsilon_r<1$. This proves the last assertion, including $r=0$,
where the first-step entropy is zero.
\end{proof}
\lean{SU2PointEntropy}{forward_entropy_finite; entropy_geometric_after_first}
\lean{SU2PointEntropy}{point_start_entropy_finite; entropy_tendsto_zero_after_first}
The explicit upper bound $4|r|$ is a printed human consequence of the mapped
domination estimate; the source states finiteness and the exact geometric
estimate with the actual first-step entropy. No sharpness of the numerical
constant is claimed. The rate here is fixed-coupling: the displayed
$\varepsilon_r$ tends to zero as $|r|\to\infty$, and this proof supplies
no coupling-uniform entropy rate. Neither it nor the boundary obstruction
addresses volume or continuum uniformity.


\section{A two-step Haar reset through a zero-trace band}
The replacement kernel has a more informative positive-coupling reset
than its global density bound alone suggests. Keep the \emph{same}
$P_r$ from the preceding section, and assume $r\geq0$. For $\delta>0$ put
\[
 B_\delta=\{Z\in G:|h(Z)|<\delta\},\qquad
 a_\delta=m(B_\delta),\qquad
 b_{r,\delta}=a_\delta e^{-(1+3\delta)r}.
\]
The set is open since $h$ is continuous. It is nonempty, since
$\operatorname{diag}(i,-i)\in\SU$ has zero trace. Thus $a_\delta>0$.
For completeness, Haar probability on a compact group assigns positive
mass to every nonempty open set: its left translates cover the group,
a finite subcover exists by compactness, and zero mass of the original
set would force total mass zero. Hence no unproved quantitative band
mass estimate is needed here.

\begin{theorem}[Actual two-step minorization]\label{thm:bandreset}
For every $C\in G$,
\begin{equation}\label{eq:bandreset}
 P_r^2(C,\cdot)\geq b_{r,\delta}m,\qquad
 0<b_{r,\delta}\leq1.
\end{equation}
The composition is of the actual probability kernel, twice; it is not
a proposed transition density or a stationary-reference reset.
\end{theorem}
\begin{proof}
For $Z\in B_\delta$, the first density obeys
\[
 k_r(C,Z)\geq e^{-(1+\delta)r},
\]
because $rh(C)h(Z)\geq-r\delta$ and $z_r(C)\leq e^r$.
At the second step, $z_r(Z)\leq e^{r\delta}$ and
$rh(Z)h(D)\geq-r\delta$ for every $D$. Thus
\[
 k_r(Z,D)\geq e^{-2r\delta}\qquad(Z\in B_\delta).
\]
For any Borel $A$, positivity and Tonelli give
\begin{align*}
 P_r^2(C,A)
 &=\int_G k_r(C,Z)\!\int_A k_r(Z,D)m(dD)m(dZ)\\
 &\geq\int_{B_\delta}e^{-(1+\delta)r}
                    e^{-2r\delta}m(A)m(dZ)
 =b_{r,\delta}m(A).
\end{align*}
The lower coefficient is positive. Evaluating the inequality at $A=G$
gives $b_{r,\delta}\leq1$, using the Markov property of the composition.
\end{proof}
\lean{SU2ZeroTraceBridge}{band_mass_pos; density_into_band; density_from_band}
\lean{SU2ZeroTraceBridge}{two_step_reset; bridge_pos; bridge_le_one}

\subsection*{Iteration, stationary convergence and the pair count}
Put $\alpha=b_{r,\delta}/2\in(0,1)$ and define the genuine residual
\[
 L(C,\cdot)=\frac{P_r^2(C,\cdot)-\alpha m}{1-\alpha}.
\]
Nonnegativity, measurable evaluations and mass one follow just as in
the previous residual construction. Here $m$ is a common reset law;
it is not asserted stationary. In fact
$\pi_r=\alpha m+(1-\alpha)\pi_r L$ follows from stationarity of $\pi_r$
for $P_r^2$, and does not give $\pi_r L=\pi_r$.
For a bounded real measurable $f$, let
$\operatorname{osc}(f)=\sup_C f(C)-\inf_C f(C)$. Since a Markov average
stays between the infimum and supremum of its input,
\[
 \operatorname{osc}(P_r^2f)
 =(1-\alpha)\operatorname{osc}(Lf)
 \leq(1-\alpha)\operatorname{osc}(f).
\]
Induction therefore proves, for $0\leq f\leq1$,
\begin{equation}\label{eq:bandosc}
 |P_r^{2n}f(C)-P_r^{2n}f(D)|
 \leq(1-\alpha)^n\leq e^{-n\alpha}.
\end{equation}
All powers are actual kernel powers. Measurability and boundedness
are preserved by kernel integration, so every induction step is defined.
The same estimate holds for $2n+1$ steps by applying the even-step bound
to $P_rf$, which still takes values in $[0,1]$.
\lean{SU2ZeroTraceBridge}{twoAverage_eq_transfer_twice; iterate_range}
\lean{SU2ZeroTraceBridge}{iterate_oscillation; iterate_exponential_bound}

A useful human consequence is convergence from every initial probability
$\mu$ to the already constructed stationary $\pi_r$. Define
$d_{\rm TV}(\mu,\nu)=\sup_A|\mu(A)-\nu(A)|$. Integrating the
two-state bound for $f=1_A$ against $\mu(dC)\pi_r(dD)$ and using
stationarity gives
\[
 d_{\rm TV}(\mu P_r^{2n},\pi_r),\quad
 d_{\rm TV}(\mu P_r^{2n+1},\pi_r)
 \leq(1-\alpha)^n.
\]
This also proves uniqueness: if $\nu$ is stationary, substitute $\nu$
and let $n\to\infty$. This total-variation conclusion is a printed
human deduction, not a new mapped formal theorem.

For a prescribed $t\geq0$, a sufficient integer pair count is
\begin{equation}\label{eq:bandcount}
 n\geq\frac{2t}{a_\delta}e^{(1+3\delta)r}.
\end{equation}
Multiplication by $\alpha=(a_\delta/2)e^{-(1+3\delta)r}$ gives
$n\alpha\geq t$, and \eqref{eq:bandosc} then gives $e^{-t}$.
These are $n$ pairs, hence $2n$ original transitions. For each
$\eta>0$, choose $\delta=\eta/3$; the positive two-step Haar coefficient
is then $a_{\eta/3}e^{-(1+\eta)r}$. Its prefactor depends on $\eta$.
No exponent-one endpoint, optimality or coupling-uniform rate is claimed.
\lean{SU2ZeroTraceBridge}{sufficient_pair_count; reset_rate_near_one}
The oscillation/total-variation mechanism is classical Doeblin contraction;
Gaubert--Qu~\cite{gq}, introductory equations (1)--(4) and Theorem 1.1,
give its coefficient and oscillation context. No new general ergodicity
principle is asserted. This Haar minorization is also not a license to
reuse Lemma~\ref{lem:residual} with reference $\pi_r$: that entropy proof
required a \emph{stationary-reference} minorization and invariant residual.
The earlier one-step entropy conclusion retains its separate proof.

\subsection*{The zero-trace row distinguishes Wilson transport}
There is a direct obstruction to identifying this reset mechanism with
one-word Wilson transport. Define the latter on the same group by
\[
 W_r(C,dD)=\frac{e^{r h(C^{-1}D)}}{\zeta_r}m(dD),\qquad
 \zeta_r=\int_Ge^{rh(V)}m(dV).
\]
Left invariance makes the denominator independent of $C$. Equivalently,
$D=CV$, where $V$ has the Haar tilt $e^{rh(V)}/\zeta_r$.
This is the source's one-word transition evaluated at the inverse
boundary argument. Its normalized relative-holonomy mean is
\[
 \int_Gh(C^{-1}D)W_r(C,dD)
 =c_r,\qquad c_r=\frac{\int_G h(V)e^{rh(V)}m(dV)}{\zeta_r}.
\]
In contrast, $h(C)=0$ makes $P_r(C,\cdot)=m$ exactly, since its
interaction vanishes and its partition is one.
\lean{SU2ZeroTraceBridge}{zero_trace_row; wilson_relative_mean}

Here is a direct human strictness argument, valid at every $r>0$.
The central transformation $V\mapsto -V$ preserves Haar probability
and sends $h$ to $-h$. Averaging the original numerator with its
transformed expression gives
\[
 \int h(V)e^{rh(V)}m(dV)
 =\int h(V)\sinh(rh(V))m(dV)>0.
\]
The integrand is nonnegative and continuous, and strictly positive
on a nonempty open neighborhood of the identity. The same compact
Haar support argument gives the strict integral inequality. All
integrands are bounded at each fixed $r$. Thus $c_r>0$, whereas
$\int h(C^{-1}D)m(dD)=\int h\,dm=0$ by the same sign symmetry.
At any zero-trace $C$ the two rows consequently differ for every $r>0$.
The mapped formal mismatch is stated eventually at large coupling,
using $c_r\to1$; the all-positive-coupling strengthening here is a
labelled human symmetry proof, not that literal formal quantifier.
\lean{SU2ZeroTraceBridge}{zero_trace_wilson_mismatch}

This distinction already occurs for one group variable. It does not
settle a Wilson lattice gap or prohibit another valid minorization
for $W_r$. In particular the algebraic rescaling
$2ne^{-(1+3\delta)r}$ in the loop oscillation exponent is only a
proposed distance scale until related independently to an actual
lattice spacing. The source labels that scale as proposed; no physical
mass gap follows by renaming a transition count.

\section{Interpretation and comparison with heat-bath mixing}
Conditional expectation and the entropy chain rule are classical. This
paper makes no novelty claim for the obstruction that an exactly retained
variable can preserve information. Raginsky's primary treatment of strong
data processing \cite{raginsky}, Definition~1, formulates strict entropy
contraction as a ratio bounded below one. The stationary tests above
exhibit ratio one in the specified channel. His finite-alphabet results
are background for the terminology; we do not substitute them for the
continuous-boundary proof.

Caputo--Menz--Tetali \cite{cmt}, Section~1.1 and Proposition~1.1, treat
heat-bath dynamics that update all coordinate positions, with approximate
entropy tensorization supplying global control. That is a different update
from repeating one boundary-preserving conditional step. Their weak
dependence hypotheses are not established here, and no such mixing estimate
is imported. Our stationary family does not contradict entropy decay for
a dynamics whose allowed updates also change the boundary.

There is also no conflict with strict mixing of a finite transfer kernel
on a different state space. The finite seven-link transfer construction in
the companion paper updates its own chain state; it is not identified here
with $T_\beta$. A kernel preserving $U$ exactly is structurally incapable
of forgetting an arbitrary initial distribution of $U$, regardless of how
positive its conditional link density is. Boundary reweighting in
\eqref{eq:actualfactor} specifies which global reference is at issue.

For the restricted class $J_{\mathsf P}=\tau_\beta$, formula
\eqref{eq:update} sends $J$ to $Q_\beta$ in one step. For arbitrary boundary
inputs the global fraction fails at every fixed finite coupling, rather than
only in a large-coupling asymptotic. These two exact statements delimit the
scope without claiming an entropy inequality that the update cannot satisfy.

Finally the hemisphere event uses a particular matrix coordinate and is
not asserted gauge invariant. The theorem concerns the stated actual
probability space and conditional update; it does not classify a gauge
quotient, all lattice block dynamics, infinite-volume measures or continuum
Schwinger functions. No lattice-spacing or volume limit is constructed, and
no physical mass-gap implication, positive or negative, follows.

\begin{thebibliography}{9}
\bibitem{raginsky} M. Raginsky, \emph{Logarithmic Sobolev inequalities and
strong data processing theorems for discrete channels}, IEEE International
Symposium on Information Theory (2013).
\href{https://maxim.ece.illinois.edu/pubs/raginsky_ISIT13.pdf}{Primary author PDF}.
\bibitem{cmt} P. Caputo, G. Menz and P. Tetali,
\emph{Approximate tensorization of entropy at high temperature},
Ann. Fac. Sci. Toulouse Math. (6) \textbf{24} (2015), 691--716.
\href{https://www.numdam.org/item/AFST_2015_6_24_4_691_0.pdf}{Primary journal PDF}.
\bibitem{gq} S. Gaubert and Z. Qu,
\emph{Dobrushin's ergodicity coefficient for Markov operators on cones},
primary author manuscript, Section 1, equations (1)--(4), Theorem 1.1.
\url{https://www.cmap.polytechnique.fr/~gaubert/PAPERS/GaubertQuIEOTD14QuFinal.pdf}.
\end{thebibliography}
\end{document}
